\subsection*{Parte 4: ¡Adivina qué sale!}
\begin{center}
\renewcommand{\arraystretch}{1.3}
\begin{tabular}{l l P{7.5cm}}
\toprule
\textbf{Código} & \textbf{Muestra} & \textbf{Por qué} \\
\midrule
\py|print(9 / 3)| & \texttt{3.0} & Dividir siempre da un decimal. \\
\py|print("4" + "4")| & \texttt{44} & Son strings, así que \texttt{+} los une. \\
\py|print("Perro" == "perro")| & \texttt{False} & En Python, las mayúsculas cuentan. \\
\py|print("pizza"[1])| & \texttt{i} & Python cuenta desde 0, así que \texttt{[1]} es la segunda letra. \\
\py|print(round(4.5))| & \texttt{4} & Un .5 va al número par más cercano. \\
\py|print(17 % 5)| & \texttt{2} & $17 = 3 \times 5 + 2$, así que sobran 2. \\
\bottomrule
\end{tabular}
\end{center}

\subsection*{Parte 4, Misión 1: Lanzador de dados}
Para 100 tiradas, cambia la línea 3 por \py|for i in range(100):| y la línea 8 por \py|str(total / 100)|. Con dos dados, el total más grande es $6 + 6 = 12$.

\subsection*{Parte 4, Misión 2: Impresora de tablas}
Para la tabla del 9, cambia la línea 1 por \py|tabla = 9|. Para llegar hasta 20, cambia la línea 3 por \py|range(20)|. Para los números al cuadrado, la línea 4 queda \py|print(f"{n} x {n} = {n * n}")|. El \texttt{True} o \texttt{False} del final te dice si el resultado es par.

\subsection*{Parte 4, Misión 3: Arte de puntos al azar}
Con 7 colores, la línea 15 necesita \py|random.randint(0, 6)|. Para la mitad de arriba, la línea 12 queda \py|y = random.randint(0, 180)|.

\subsection*{Parte 4, Misión 4: Máquina de nombres}
Dice \py|i + 1| porque \py|i| empieza en 0, pero las personas cuentan las letras desde 1.

\subsection*{Parte 4: Tradúcelo tú}
Tu código puede verse un poco diferente y aun así estar bien. Si imprime tres totales, ¡lo lograste!
\begin{codenonum}
import random
for i in range(3):
    a = random.randint(1, 10)
    b = random.randint(1, 10)
    print("Total: " + str(a + b))
\end{codenonum}
Una ejecución real imprimió \texttt{Total: 9}, \texttt{Total: 19} y \texttt{Total: 8}. Un f-string también funciona: \py|print(f"Total: {a + b}")|.

\subsection*{Parte 4: Acertijo de dos líneas}
\begin{codenonum}
palabra = "Tortuga"
print(palabra[0] + str(len(palabra)))
\end{codenonum}
\py|palabra[0]| es la primera letra, \texttt{T}. \py|len(palabra)| es 7, y \py|str( )| lo convierte en string para poder unirlo.
